\section{Second derivatives and singular integrals}
\label{sec:hessian-singular-integrals}

Let $d\geq1$. Write $D^2v=(\partial_a\partial_bv)_{a,b=1}^d$, and use Euclidean norms for vectors and the Hilbert--Schmidt norm $|M|=(\sum_{a,b}|M_{ab}|^2)^{1/2}$ for matrices in $L^p$ norms. Constants $C_d$ and $C_{d,p}$ may change from line to line and depend only on the indicated parameters. A bound for a scalar Laplacian controls the full real Hessian in $L^p$, although it need not control that Hessian pointwise. Plancherel gives the estimate for $p=2$. The extension to other finite exponents uses cancellation away from the support of a function.

\begin{lemma}\label{lem:hessian-multiplier-kernel}
For $1\leq a,b\leq d$, define $T_{ab}$ on $L^2(\mathbb R^d)$ by
\[
 \widehat{T_{ab}f}(\xi)=-\frac{\xi_a\xi_b}{|\xi|^2}\widehat f(\xi),
 \qquad \xi\neq0,
\]
with multiplier zero at the origin. This is a self-adjoint operator of norm at most one. There is a smooth real kernel $K_{ab}$ away from zero such that
\[
 |K_{ab}(z)|\leq C_d|z|^{-d},\qquad
 |\nabla K_{ab}(z)|\leq C_d|z|^{-d-1}.
\]
If $f\in L^1\cap L^2$ has compact support, then almost everywhere outside that support
\[
 T_{ab}f(x)=\int K_{ab}(x-y)f(y)\,\mathrm d y.
\]
For $v\in C_c^\infty$, one has $\partial_a\partial_bv=T_{ab}(-\Delta v)$.
\end{lemma}
\begin{proof}
The multiplier is real and bounded by one in absolute value, so Plancherel gives the norm and adjoint assertions. The derivative identity follows by multiplying the transform of $-\Delta v$ by the displayed multiplier; both sides are $L^2$ functions.

\Needspace{16\baselineskip}
Put $H_t(z)=(4\pi t)^{-d/2}\mathrm e^{-|z|^2/(4t)}$. The Gaussian formula gives $\widehat H_t(\xi)=\mathrm e^{-4\pi^2t|\xi|^2}$. For $0<\varepsilon<R<\infty$, the operators
\[
 \int_\varepsilon^R\partial_a\partial_bH_t*f\,\mathrm d t
\]
have multiplier
\[
 -\frac{\xi_a\xi_b}{|\xi|^2}
    \bigl(\mathrm e^{-4\pi^2\varepsilon|\xi|^2}
          -\mathrm e^{-4\pi^2R|\xi|^2}\bigr)
\]
and converge to $T_{ab}f$ in $L^2$ as $\varepsilon\downarrow0$, $R\uparrow\infty$. Each finite integral is justified by the integrable smooth convolution kernels; bounded frequency multipliers and dominated convergence give the limit.

For $z\neq0$ define
\[
 K_{ab}(z)=\int_0^\infty\partial_a\partial_bH_t(z)\,\mathrm d t.
\]
Every differentiated integrand is bounded by a polynomial in $|z|/\sqrt t$ times $C t^{-(d+2+l)/2}\mathrm e^{-|z|^2/(4t)}$ for $l$ further derivatives. This is integrable uniformly on compact sets away from zero: the exponential controls $t\downarrow0$, and $t^{-(d+2+l)/2}$ controls $t\uparrow\infty$. Thus the kernel is smooth there. The change $t=r^2s$ shows $K_{ab}(rz)=r^{-d}K_{ab}(z)$ for $r>0$. Its bounded derivatives on the unit sphere give the two estimates.

On a compact set disjoint from the support of $f$, the integral of $|\partial_a\partial_bH_t(x-y)|$ in $t$ is at most $C_d|x-y|^{-d}$. Fubini and dominated convergence therefore identify the finite heat integrals with the asserted convolution in the limit, uniformly on that compact set. Their $L^2$ limit is $T_{ab}f$, proving the almost-everywhere identity. Exhaust the complement by such compact sets. This argument also applies for $d=1$, where $T_{11}=-1$ and the off-origin kernel vanishes.
\end{proof}

\begin{lemma}\label{lem:hessian-multiplier-weak-one}
For $f\in L^1(\mathbb R^d)\cap L^2(\mathbb R^d)$ and $\lambda>0$,
\[
 |\{x:|T_{ab}f(x)|>\lambda\}|\leq\frac{C_d}{\lambda}\|f\|_1.
\]
\end{lemma}
\begin{proof}
The dyadic cubes are $2^k(m+[0,1)^d)$, for $k\in\mathbb Z$ and $m\in\mathbb Z^d$. Choose the maximal cubes $Q$ for which the average of $|f|$ exceeds $\lambda$. Such a cube lies below a maximal one: the averages on its successive ancestors tend to zero, since their volumes tend to infinity and $f\in L^1$. Two maximal cubes are disjoint, and their parent averages give
\[
 \sum_Q|Q|\leq\lambda^{-1}\|f\|_1,\qquad
 \lambda<\frac1{|Q|}\int_Q|f|\leq2^d\lambda.
\]
Outside their union, $|f|\leq\lambda$ almost everywhere. Indeed a shrinking dyadic cube containing a Lebesgue point is contained in a centered cube of comparable volume; \cref{thm:lebesgue-differentiation} makes its average tend to $|f(x)|$.

Write $f_Q=|Q|^{-1}\int_Qf$, set $g=f$ off the selected cubes and $g=f_Q$ on each of them, and put $b_Q=(f-f_Q)\mathbf1_Q$. Then
\[
 \|g\|_1\leq\|f\|_1,\quad \|g\|_\infty\leq2^d\lambda,
 \quad \|g\|_2^2\leq2^d\lambda\|f\|_1,
 \qquad \sum_Q\|b_Q\|_1\leq2\|f\|_1.
\]
Each $b_Q$ has integral zero. Jensen's inequality and disjoint supports give $g\in L^2$ and convergence of $\sum_Qb_Q=f-g$ in $L^2$.

Let $c_Q$ be the center of $Q$ and let $Q^*$ be its concentric enlargement by the factor $4\sqrt d$. For $x\notin Q^*$, the segment from $x-c_Q$ to $x-y$, $y\in Q$, stays at distance at least $|x-c_Q|/2$ from zero. Cancellation and \cref{lem:hessian-multiplier-kernel} yield
\[
 \begin{aligned}
 |T_{ab}b_Q(x)|
 &=\left|\int_Q\bigl(K_{ab}(x-y)-K_{ab}(x-c_Q)\bigr)b_Q(y)\,\mathrm d y\right|\\
 &\leq C_d\frac{\operatorname{diam}Q}{|x-c_Q|^{d+1}}\|b_Q\|_1.
 \end{aligned}
\]
Polar integration outside a ball of radius comparable to $\operatorname{diam}Q$ gives
\[
 \int_{\mathbb R^d\setminus Q^*}|T_{ab}b_Q|\leq C_d\|b_Q\|_1.
\]
Put $B^*=\bigcup_QQ^*$. Its volume is at most $C_d\lambda^{-1}\|f\|_1$. Finite partial sums of the $b_Q$ converge in $L^2$, and so do their images under $T_{ab}$. Select an almost-everywhere convergent subsequence and use Fatou outside $B^*$ to get
\[
 \int_{\mathbb R^d\setminus B^*}|T_{ab}(f-g)|\leq C_d\|f\|_1.
\]
Chebyshev's inequality for this integral and the $L^2$ bound for $T_{ab}g$ bound the sets where the two summands exceed $\lambda/2$. Including $B^*$ proves the assertion.
\end{proof}
\index{Calder\'on--Zygmund decomposition}
\index{singular integral!weak $(1,1)$ estimate}

\begin{theorem}\label{thm:hessian-multiplier-lp}
For every $1<p<\infty$, $T_{ab}$ extends uniquely to a bounded operator on $L^p(\mathbb R^d)$. In particular
\begin{equation}\label{eq:euclidean-hessian-lp}
 \|D^2v\|_p\leq C_{d,p}\|\Delta v\|_p,
 \qquad v\in C_c^\infty(\mathbb R^d).
\end{equation}
\end{theorem}
\begin{proof}
Let $1<p<2$ and $f\in L^p\cap L^2$. For $\lambda>0$ split $f$ into $f\mathbf1_{\{|f|>\lambda\}}$ and $f\mathbf1_{\{|f|\leq\lambda\}}$. The first lies in $L^1\cap L^2$. Apply the weak estimate to it, and the $L^2$ estimate to the second, both with threshold $\lambda/2$. This gives
\[
 |\{|T_{ab}f|>\lambda\}|
 \leq\frac{C_d}{\lambda}\int_{|f|>\lambda}|f|
      +\frac{C_d}{\lambda^2}\int_{|f|\leq\lambda}|f|^2.
\]
For any measurable $h$, Tonelli applied to $|h(x)|^p=\int_0^{|h(x)|}p\lambda^{p-1}\,\mathrm d\lambda$ gives its norm as the integral of this distribution function, including infinite norms. Multiplying the inequality by $p\lambda^{p-1}$ and integrating therefore gives
\[
 \|T_{ab}f\|_p^p
 \leq C_dp\left(\frac1{p-1}+\frac1{2-p}\right)\|f\|_p^p.
\]
The two inner integrals are $\int_0^{|f|}\lambda^{p-2}\,\mathrm d\lambda$ and $\int_{|f|}^\infty\lambda^{p-3}\,\mathrm d\lambda$. Their stated values explain both restrictions on $p$. Plancherel handles $p=2$.

For $p>2$ put $p'=p/(p-1)<2$, take $f\in L^p\cap L^2$, and write $h=T_{ab}f$. For integers $N\geq1$, on $E_N=B(0,N)\cap\{|h|\leq N\}$ set $g_N=\mathbf1_{E_N}h|h|^{p-2}$. This belongs to $L^2\cap L^{p'}$. Self-adjointness, the proved $p'$ estimate and H\"older give
\[
 \int_{E_N}|h|^p=\int h\overline{g_N}
 =\int f\overline{T_{ab}g_N}
 \leq C_{d,p'}\|f\|_p\left(\int_{E_N}|h|^p\right)^{1/p'}.
\]
If the integral is nonzero, division bounds its $p$th root; if it is zero the bound holds directly. The sets $E_N$ increase to full measure, since $h\in L^2$. Monotone convergence proves the $p$ estimate. Density of $C_c^\infty$ in $L^p$, from \cref{lem:mollifier-density}, gives the unique extension. Finally apply the derivative identity in \cref{lem:hessian-multiplier-kernel} to each coefficient of $D^2v$.
\end{proof}
\index{singular integral!$L^p$ boundedness}
\index{Hessian!$L^p$ estimate}

For an integer $k\geq0$, $1\leq p<\infty$ and an open set $U\subset\mathbb R^d$, the Sobolev space\index{Sobolev space} $W^{k,p}(U)$ consists of $L^p$ functions whose distributional derivatives of order at most $k$ belong to $L^p$. Its norm is the sum of those derivative norms. On a compact smooth manifold, a finite coordinate cover and a smooth partition of unity define equivalent such norms. The chain rule and bounded derivatives of the coordinate changes prove equivalence for smooth functions; local smoothing and distributional passage to the limit give the same rule for weak derivatives.

\begin{lemma}\label{lem:gradient-lp-interpolation}
For $1\leq p<\infty$ and $\varepsilon>0$ there is $C_{d,\varepsilon}$ such that
\[
 \|\nabla v\|_p\leq\varepsilon\|D^2v\|_p+C_{d,\varepsilon}\|v\|_p,
 \qquad v\in C_c^\infty(\mathbb R^d).
\]
\end{lemma}
\begin{proof}
Taylor's formula along the $a$th coordinate direction gives, for $h>0$,
\[
 h\partial_av(x)=v(x+he_a)-v(x)
       -\int_0^h(h-t)\partial_a^2v(x+te_a)\,\mathrm d t.
\]
Translation invariance and the integral triangle inequality give $\|\partial_av\|_p\leq2h^{-1}\|v\|_p+(h/2)\|\partial_a^2v\|_p$. Sum over $a$ and choose $h$ sufficiently small in terms of $d,\varepsilon$.
\end{proof}

For $0<\alpha\leq1$, the norm $\|v\|_{C^{0,\alpha}(F)}$ on a compact set $F$ is $\sup_F|v|+\sup_{x\neq y\in F}|v(x)-v(y)|/|x-y|^\alpha$. The notation $C^{1,\alpha}$ requires continuously differentiable functions with first derivatives H\"older continuous with exponent $\alpha$.

\Needspace{10\baselineskip}
\begin{lemma}[Morrey's inequality]\label{lem:morrey-sobolev-representative}
Let $d<p<\infty$, $\alpha=1-d/p$, and $B'\Subset B$ be Euclidean balls. Every $v\in W^{1,p}(B)$ has a continuous representative with
\[
 \|v\|_{C^{0,\alpha}(\overline{B'})}\leq C\|v\|_{W^{1,p}(B)}.
\]
Here $C$ depends only on $d,p,B',B$. Consequently a $W^{2,p}_{\mathrm{loc}}$ function has a $C^{1,\alpha}_{\mathrm{loc}}$ representative.
\end{lemma}
\begin{proof}
First let $v$ be smooth. If $Q$ is a ball of radius $r$ and $x\in Q$, integrate $v(x)-v(z)$ along the segment from $x$ to $z$, then average in $z\in Q$. Polar coordinates about $x$ give
\[
 |v(x)-v_Q|\leq C_d\int_Q|\nabla v(z)|\,|x-z|^{1-d}\,\mathrm d z,
 \qquad v_Q=\frac1{|Q|}\int_Qv.
\]
To see the constant, write $z=x+s\theta$, whose allowed radii satisfy $0<s<R(\theta)\leq2r$ by convexity. Interchanging the segment integral and the radial integral produces the factor $\int_t^{R(\theta)}s^{d-1}\,\mathrm d s/|Q|\leq C_d$; converting back to volume gives $t^{1-d}$. H\"older now bounds the displayed integral by $C_{d,p}r^{1-d/p}\|\nabla v\|_{L^p(Q)}$. Indeed $(1-d)p'>-d$ exactly when $p>d$; the same calculation holds for $d=1$.

For $x,y\in B'$ with $r=|x-y|$ sufficiently small compared with $\operatorname{dist}(B',\partial B)$, use the common ball $Q=B(x,2r)\subset B$. Applying this bound at both points gives
\[
 |v(x)-v(y)|\leq C_{d,p}|x-y|^{1-d/p}\|\nabla v\|_{L^p(B)}.
\]
Using a fixed small ball about each $x\in B'$ also bounds $|v(x)|$ by $C\|v\|_{W^{1,p}(B)}$, since its average is bounded by $|Q|^{-1/p}\|v\|_p$. For more distant pairs this supremum bound gives the same H\"older estimate with a constant depending on $B',B$.

For a general $W^{1,p}$ function, mollify on a slightly smaller ball still containing $\overline{B'}$. The functions and their weak derivatives converge in $L^p$ by \cref{lem:mollifier-density,prop:distribution-smoothing}. The just proved estimate makes the approximations Cauchy uniformly and in the H\"older seminorm on $B'$. Their continuous limit agrees almost everywhere with $v$ and obeys the estimate. For $v\in W^{2,p}$ apply this argument to $v$ and each first derivative. Their smooth approximations converge locally uniformly; passing to the limit in the segment formula identifies the continuous derivative limits as the classical derivatives of the continuous limit. This gives the last assertion.
\end{proof}
\index{Morrey's inequality}
\index{Sobolev space!continuous representatives}
